% ====================================================================
%  Ablations section for the Predictive Evaluation Challenge report.
%  Drop this file into Overleaf and `% ====================================================================
%  Ablations section for the Predictive Evaluation Challenge report.
%  Drop this file into Overleaf and `% ====================================================================
%  Ablations section for the Predictive Evaluation Challenge report.
%  Drop this file into Overleaf and `% ====================================================================
%  Ablations section for the Predictive Evaluation Challenge report.
%  Drop this file into Overleaf and `\input{ablations_section.tex}` from
%  the main .tex, OR copy the body wholesale into the relevant section.
%
%  Required preamble packages (add to main.tex if not already there):
%      \usepackage{booktabs}
%      \usepackage{makecell}
%      \usepackage{array}
%      \usepackage{multirow}
%      \usepackage{xcolor}
%      \usepackage{amsmath}
%
%  Score macros — once the leaderboard returns each round, replace the
%  TBD with the number. Negative log-loss is signed so a sign is
%  required (e.g. -0.65, NOT 0.65).
% ====================================================================
\newcommand{\TBD}{\textit{pending}}
% -- known scores from prior rounds (sub 1..8) -----------------------
\newcommand{\lbNLLsubOne}{-1.01}\newcommand{\lbAUCsubOne}{0.62}
\newcommand{\lbNLLsubTwo}{-0.93}\newcommand{\lbAUCsubTwo}{0.63}
\newcommand{\lbNLLsubThree}{-0.85}\newcommand{\lbAUCsubThree}{0.63}
\newcommand{\lbNLLsubFive}{-0.65}\newcommand{\lbAUCsubFive}{0.63}
\newcommand{\lbNLLsubSix}{-0.65}\newcommand{\lbAUCsubSix}{\TBD}
\newcommand{\lbNLLsubSeven}{-0.65}\newcommand{\lbAUCsubSeven}{\TBD}
\newcommand{\lbNLLsubEight}{-0.75}\newcommand{\lbAUCsubEight}{\TBD}
% -- Part-V ablations: fill in as rounds resolve ---------------------
\newcommand{\lbNLLconstHalf}{\TBD}\newcommand{\lbAUCconstHalf}{\TBD}
\newcommand{\lbNLLconstGM}{\TBD}\newcommand{\lbAUCconstGM}{\TBD}
\newcommand{\lbNLLsubjMean}{\TBD}\newcommand{\lbAUCsubjMean}{\TBD}
\newcommand{\lbNLLbmCond}{\TBD}\newcommand{\lbAUCbmCond}{\TBD}
\newcommand{\lbNLLblendSubj}{\TBD}\newcommand{\lbAUCblendSubj}{\TBD}
\newcommand{\lbNLLlabeledShift}{\TBD}\newcommand{\lbAUClabeledShift}{\TBD}
% ====================================================================

\section{Ablations and Strategy Coverage}
\label{sec:ablations}

We organise our exploration around Discussion~6's two scaffolds:
(i)~the seven-rung calibration climb (from a constant 0.5 up to a
content-aware PGE with adaptive labels), and
(ii)~the six-family option space
(F1 calibrated priors, F2 IRT alone, F3 content-based, F4 metadata
stabilisers, F5 local LLM judges, F6 adaptive labels).
Table~\ref{tab:coverage} maps each submission we ran onto these axes,
and Table~\ref{tab:ablations} reports the leaderboard outcomes.

\subsection{Coverage map}

\begin{table}[h]
\centering
\small
\caption{Coverage of Discussion 6 strategy axes. Rows are the
seven-rung calibration climb; columns are the six families.
\ding{51} marks an ablation we shipped to the leaderboard; \ding{55} marks
families we deliberately did not pursue (with justification in
\S\ref{sec:ablations-skipped}).}
\label{tab:coverage}
\begin{tabular}{lcccccc}
\toprule
                                  & \textbf{F1} & \textbf{F2} & \textbf{F3} & \textbf{F4} & \textbf{F5} & \textbf{F6} \\
                                  & priors      & IRT         & content     & metadata    & local LLM   & adaptive \\
\midrule
R1: $\hat p = 0.5$                & \checkmark  &             &             &             &             &             \\
R2: global mean                   & \checkmark  &             &             &             &             &             \\
R3: smoothed subject mean         & \checkmark  &             &             & \checkmark  &             &             \\
R4: per-(benchmark, condition) mean & \checkmark &             &             & \checkmark  &             &             \\
R5: + content (text $\to$ params) &             &             & \checkmark  &             & \ding{55}   &             \\
R6: + IRT $\theta$ (full PGE)     &             & \ding{55}   & \checkmark  &             &             &             \\
R7: + adaptive K{=}5 labels       &             &             & \checkmark  &             &             & \checkmark  \\
\bottomrule
\end{tabular}
\end{table}

The lower four rungs are pure ablations of slide~30's recipe; the
upper three are variants of our winning submission (Sub~5, T-scaled
mpnet) along the metadata and adaptive-label axes.

\subsection{Results}

\begin{table*}[ht]
\centering
\small
\caption{All submissions, ordered by leaderboard NLL (negative
log-loss; higher is better). LB columns are filled in as each round
resolves. Rungs and families reference Discussion~6 Part~V
slides~29 and~31.}
\label{tab:ablations}
\renewcommand{\arraystretch}{1.15}
\begin{tabular}{rl c c l rr}
\toprule
\textbf{ID} & \textbf{Submission} & \textbf{Rung} & \textbf{Family} & \textbf{Mechanism} & \textbf{LB NLL} & \textbf{LB AUC} \\
\midrule
\multicolumn{7}{l}{\textit{Climb baselines (priors / metadata)}}\\
9  & \texttt{const\_0p5}       & R1 & F1     & constant $0.5$                                       & \lbNLLconstHalf      & \lbAUCconstHalf \\
10 & \texttt{const\_glob\_mean} & R2 & F1     & constant $\bar y_{\text{train}}{=}0.645$             & \lbNLLconstGM        & \lbAUCconstGM   \\
15 & \texttt{subj\_mean\_lookup}& R3 & F1+F4  & $0.9\,\bar y_s + 0.1\,\bar y_{\text{train}}$         & \lbNLLsubjMean       & \lbAUCsubjMean  \\
16 & \texttt{bm\_cond\_lookup} & R4 & F1+F4  & $\bar y_{(b,c)}$, fallback $\bar y_{\text{train}}$   & \lbNLLbmCond         & \lbAUCbmCond    \\
\midrule
\multicolumn{7}{l}{\textit{Prior context: encoder upgrade alone vs. with T-scaling}}\\
1 & \texttt{mpnet\_platt\_meanresid}     & R6 & F3 & mpnet+Platt+per-subject offset  & \lbNLLsubOne   & \lbAUCsubOne   \\
2 & \texttt{mpnet\_noplatt\_meanresid}   & R6 & F3 & sub 1 -- Platt                  & \lbNLLsubTwo   & \lbAUCsubTwo   \\
3 & \texttt{mpnet\_no\_offset}           & R6 & F3 & sub 2 -- offset                 & \lbNLLsubThree & \lbAUCsubThree \\
\textbf{5} & \textbf{\texttt{tscale\_mpnet} (winner)} & R6 & F3 & \textbf{sub 3 + post-hoc T-scaling ($T{=}4.07$)} & \textbf{\lbNLLsubFive} & \textbf{\lbAUCsubFive} \\
6  & \texttt{tscale\_T3p0}              & R6 & F3 & sub 5 with $T{=}3.0$              & \lbNLLsubSix    & \lbAUCsubSix     \\
7  & \texttt{tscale\_T5p5}              & R6 & F3 & sub 5 with $T{=}5.5$              & \lbNLLsubSeven  & \lbAUCsubSeven   \\
8  & \texttt{tscale\_plus\_offset}      & R6 & F3+F4 & sub 5 + per-subject mean-resid    & \lbNLLsubEight  & \lbAUCsubEight   \\
\midrule
\multicolumn{7}{l}{\textit{Part-V ablations on top of Sub 5}}\\
13 & \texttt{sub5\_blend\_subj}     & R6+R3 & F3+F4 & $0.5\,p_{\text{sub5}} + 0.5\,\bar y_s$       & \lbNLLblendSubj    & \lbAUCblendSubj    \\
14 & \texttt{sub5\_labeled\_shift}  & R7    & F3+F6 & sub 5 + shrunk global offset from $K{=}5$ labels, $\lambda{=}25$, $|\Delta|{\le}0.05$ & \lbNLLlabeledShift & \lbAUClabeledShift \\
\bottomrule
\end{tabular}
\end{table*}

\subsection{What each ablation tests}

\paragraph{R1: \texttt{const\_0p5} (Sub~9).}
The information-theoretic floor: $-\ln 2 \approx -0.6931$. Anything
that does not beat this is worse than no model. We include it
because the gap between $-0.6931$ and our best (Sub~5 at
\lbNLLsubFive) quantifies how much of the leaderboard score comes
from prior calibration alone.

\paragraph{R2: \texttt{const\_glob\_mean} (Sub~10).}
The Bayes-optimal constant under squared error: $\bar y =
\sum_i y_i / N = 0.645$ from the training response matrix. If the
leaderboard's marginal pass rate matches training, this submission
recovers the entire R1\,$\to$\,R2 step from slide~29 ``for free.''

\paragraph{R3: \texttt{subj\_mean\_lookup} (Sub~15).}
The smoothed-subject-mean rung. We apply a 10\,\% Bayesian shrinkage
toward the global mean, $\hat p = 0.9\,\bar y_s + 0.1\,\bar
y_{\text{train}}$, because two of 745 training subjects
(\textit{DeepSeek-V3.2-Exp-Thinking} and
\textit{gemini-2.5-flash}) have $\bar y_s = 0$ exactly and a raw
lookup would risk an unbounded NLL contribution if those subjects
appear with $y{=}1$. Predictions are clipped to $[10^{-3}, 1{-}10^{-3}]$
per slide~28.

\paragraph{R4: \texttt{bm\_cond\_lookup} (Sub~16).}
Per-benchmark mean accuracy (we collapse the condition axis because
$>$80\,\% of items use \texttt{none}). This is expected to be the
weakest of the four prior submissions: leaderboard rounds may sample
held-out benchmarks not seen in training, in which case the
predictor degenerates to the global mean. The result therefore
isolates the contribution of \emph{generalising} metadata signals
versus memorising them.

\paragraph{R6 winner: \texttt{tscale\_mpnet} (Sub~5).}
Our best-performing submission. The all-mpnet-base-v2 encoder
produces 768-d embeddings for subject, item, benchmark, and
condition text; an MLP maps the concatenated representation to
per-item IRT parameters $(a, b)$; the full PGE
$\sigma(a\,(\theta - b))$ is computed in inference. Post-hoc
temperature scaling at $T{=}4.07$ (fit on a cold-start val split via
\texttt{dump\_cs\_logits.py}) flattens the over-confident logits
produced by the cold-start mismatch.

\paragraph{Sub~5 metadata blend: \texttt{sub5\_blend\_subj} (Sub~13).}
Slide~31 lists ``F4 metadata stabilisers'' as a hedge against the
content head being uncalibrated. We blend
$\hat p = 0.5\,p_{\text{sub5}} + 0.5\,\bar y_s$. The blend is
robust: T-scaled mpnet outputs are concentrated near $0.5$, so even
when $\bar y_s = 0$ the blend lands near $0.25$ rather than
saturating. This is the variant most likely to beat the Sub~5 NLL.

\paragraph{Sub~5 with K{=}5 adaptive labels: \texttt{sub5\_labeled\_shift} (Sub~14).}
Sub~8 attempted a per-subject mean-residual offset from the $K{=}5$
labels per category and regressed by $-0.10$ NLL—five labels split
across $\approx$5 subjects yields one label per subject, which is
much too noisy. We retreat to the most conservative possible
adaptive-label lever: a \emph{single} global offset
$\Delta = \frac{1}{N+\lambda}\sum_i (y_i - \hat p_i)$ with shrinkage
$\lambda{=}25$ (so 25 labels see a $25/(25{+}25) = 50\,\%$ pull) and a
magnitude cap $|\Delta| \le 0.05$. We intentionally omit
\texttt{labeling.py} so that the platform supplies uniformly random
labels, which give an unbiased mean-residual estimator.

\subsection{Strategies we did not pursue}
\label{sec:ablations-skipped}

\paragraph{F2 IRT alone (\ding{55}).}
Slide~31 warns this ``fails cold-start'': an IRT model without text
features cannot place a previously-unseen subject in the $\theta$
space. Our winning model already uses IRT \emph{with} text features
(rung~6); a clean ``IRT-only'' ablation would require retraining
from scratch on the train matrix without text, which we judged
unlikely to be informative.

\paragraph{F5 local LLM judge (\ding{55}).}
A small generative LLM scoring items on the sandbox GPU would
require a new HuggingFace repo declared in \texttt{models.txt}, a
calibration head trained on its log-probabilities, and a re-run of
the cold-start T-scaling sweep. The sandbox imposes a 30--60~min
per-tier timeout (Table~\ref{tab:gputiers}) and the implementation
budget did not permit a credible evaluation in the time remaining.

\subsection{Headline takeaway}

The pure-prior submissions (R1--R4) bound the contribution of
\emph{calibration alone}: the R1 floor is $-0.6931$ and Sub~5 sits at
\lbNLLsubFive. The remaining gap to the leaderboard ceiling
must therefore come from the content head and/or label-conditional
shifts. \texttt{sub5\_blend\_subj} (Sub~13) and
\texttt{sub5\_labeled\_shift} (Sub~14) target exactly those two
levers in the most conservative form we could justify against Sub~8's
$-0.10$ NLL regression.
` from
%  the main .tex, OR copy the body wholesale into the relevant section.
%
%  Required preamble packages (add to main.tex if not already there):
%      \usepackage{booktabs}
%      \usepackage{makecell}
%      \usepackage{array}
%      \usepackage{multirow}
%      \usepackage{xcolor}
%      \usepackage{amsmath}
%
%  Score macros — once the leaderboard returns each round, replace the
%  TBD with the number. Negative log-loss is signed so a sign is
%  required (e.g. -0.65, NOT 0.65).
% ====================================================================
\newcommand{\TBD}{\textit{pending}}
% -- known scores from prior rounds (sub 1..8) -----------------------
\newcommand{\lbNLLsubOne}{-1.01}\newcommand{\lbAUCsubOne}{0.62}
\newcommand{\lbNLLsubTwo}{-0.93}\newcommand{\lbAUCsubTwo}{0.63}
\newcommand{\lbNLLsubThree}{-0.85}\newcommand{\lbAUCsubThree}{0.63}
\newcommand{\lbNLLsubFive}{-0.65}\newcommand{\lbAUCsubFive}{0.63}
\newcommand{\lbNLLsubSix}{-0.65}\newcommand{\lbAUCsubSix}{\TBD}
\newcommand{\lbNLLsubSeven}{-0.65}\newcommand{\lbAUCsubSeven}{\TBD}
\newcommand{\lbNLLsubEight}{-0.75}\newcommand{\lbAUCsubEight}{\TBD}
% -- Part-V ablations: fill in as rounds resolve ---------------------
\newcommand{\lbNLLconstHalf}{\TBD}\newcommand{\lbAUCconstHalf}{\TBD}
\newcommand{\lbNLLconstGM}{\TBD}\newcommand{\lbAUCconstGM}{\TBD}
\newcommand{\lbNLLsubjMean}{\TBD}\newcommand{\lbAUCsubjMean}{\TBD}
\newcommand{\lbNLLbmCond}{\TBD}\newcommand{\lbAUCbmCond}{\TBD}
\newcommand{\lbNLLblendSubj}{\TBD}\newcommand{\lbAUCblendSubj}{\TBD}
\newcommand{\lbNLLlabeledShift}{\TBD}\newcommand{\lbAUClabeledShift}{\TBD}
% ====================================================================

\section{Ablations and Strategy Coverage}
\label{sec:ablations}

We organise our exploration around Discussion~6's two scaffolds:
(i)~the seven-rung calibration climb (from a constant 0.5 up to a
content-aware PGE with adaptive labels), and
(ii)~the six-family option space
(F1 calibrated priors, F2 IRT alone, F3 content-based, F4 metadata
stabilisers, F5 local LLM judges, F6 adaptive labels).
Table~\ref{tab:coverage} maps each submission we ran onto these axes,
and Table~\ref{tab:ablations} reports the leaderboard outcomes.

\subsection{Coverage map}

\begin{table}[h]
\centering
\small
\caption{Coverage of Discussion 6 strategy axes. Rows are the
seven-rung calibration climb; columns are the six families.
\ding{51} marks an ablation we shipped to the leaderboard; \ding{55} marks
families we deliberately did not pursue (with justification in
\S\ref{sec:ablations-skipped}).}
\label{tab:coverage}
\begin{tabular}{lcccccc}
\toprule
                                  & \textbf{F1} & \textbf{F2} & \textbf{F3} & \textbf{F4} & \textbf{F5} & \textbf{F6} \\
                                  & priors      & IRT         & content     & metadata    & local LLM   & adaptive \\
\midrule
R1: $\hat p = 0.5$                & \checkmark  &             &             &             &             &             \\
R2: global mean                   & \checkmark  &             &             &             &             &             \\
R3: smoothed subject mean         & \checkmark  &             &             & \checkmark  &             &             \\
R4: per-(benchmark, condition) mean & \checkmark &             &             & \checkmark  &             &             \\
R5: + content (text $\to$ params) &             &             & \checkmark  &             & \ding{55}   &             \\
R6: + IRT $\theta$ (full PGE)     &             & \ding{55}   & \checkmark  &             &             &             \\
R7: + adaptive K{=}5 labels       &             &             & \checkmark  &             &             & \checkmark  \\
\bottomrule
\end{tabular}
\end{table}

The lower four rungs are pure ablations of slide~30's recipe; the
upper three are variants of our winning submission (Sub~5, T-scaled
mpnet) along the metadata and adaptive-label axes.

\subsection{Results}

\begin{table*}[ht]
\centering
\small
\caption{All submissions, ordered by leaderboard NLL (negative
log-loss; higher is better). LB columns are filled in as each round
resolves. Rungs and families reference Discussion~6 Part~V
slides~29 and~31.}
\label{tab:ablations}
\renewcommand{\arraystretch}{1.15}
\begin{tabular}{rl c c l rr}
\toprule
\textbf{ID} & \textbf{Submission} & \textbf{Rung} & \textbf{Family} & \textbf{Mechanism} & \textbf{LB NLL} & \textbf{LB AUC} \\
\midrule
\multicolumn{7}{l}{\textit{Climb baselines (priors / metadata)}}\\
9  & \texttt{const\_0p5}       & R1 & F1     & constant $0.5$                                       & \lbNLLconstHalf      & \lbAUCconstHalf \\
10 & \texttt{const\_glob\_mean} & R2 & F1     & constant $\bar y_{\text{train}}{=}0.645$             & \lbNLLconstGM        & \lbAUCconstGM   \\
15 & \texttt{subj\_mean\_lookup}& R3 & F1+F4  & $0.9\,\bar y_s + 0.1\,\bar y_{\text{train}}$         & \lbNLLsubjMean       & \lbAUCsubjMean  \\
16 & \texttt{bm\_cond\_lookup} & R4 & F1+F4  & $\bar y_{(b,c)}$, fallback $\bar y_{\text{train}}$   & \lbNLLbmCond         & \lbAUCbmCond    \\
\midrule
\multicolumn{7}{l}{\textit{Prior context: encoder upgrade alone vs. with T-scaling}}\\
1 & \texttt{mpnet\_platt\_meanresid}     & R6 & F3 & mpnet+Platt+per-subject offset  & \lbNLLsubOne   & \lbAUCsubOne   \\
2 & \texttt{mpnet\_noplatt\_meanresid}   & R6 & F3 & sub 1 -- Platt                  & \lbNLLsubTwo   & \lbAUCsubTwo   \\
3 & \texttt{mpnet\_no\_offset}           & R6 & F3 & sub 2 -- offset                 & \lbNLLsubThree & \lbAUCsubThree \\
\textbf{5} & \textbf{\texttt{tscale\_mpnet} (winner)} & R6 & F3 & \textbf{sub 3 + post-hoc T-scaling ($T{=}4.07$)} & \textbf{\lbNLLsubFive} & \textbf{\lbAUCsubFive} \\
6  & \texttt{tscale\_T3p0}              & R6 & F3 & sub 5 with $T{=}3.0$              & \lbNLLsubSix    & \lbAUCsubSix     \\
7  & \texttt{tscale\_T5p5}              & R6 & F3 & sub 5 with $T{=}5.5$              & \lbNLLsubSeven  & \lbAUCsubSeven   \\
8  & \texttt{tscale\_plus\_offset}      & R6 & F3+F4 & sub 5 + per-subject mean-resid    & \lbNLLsubEight  & \lbAUCsubEight   \\
\midrule
\multicolumn{7}{l}{\textit{Part-V ablations on top of Sub 5}}\\
13 & \texttt{sub5\_blend\_subj}     & R6+R3 & F3+F4 & $0.5\,p_{\text{sub5}} + 0.5\,\bar y_s$       & \lbNLLblendSubj    & \lbAUCblendSubj    \\
14 & \texttt{sub5\_labeled\_shift}  & R7    & F3+F6 & sub 5 + shrunk global offset from $K{=}5$ labels, $\lambda{=}25$, $|\Delta|{\le}0.05$ & \lbNLLlabeledShift & \lbAUClabeledShift \\
\bottomrule
\end{tabular}
\end{table*}

\subsection{What each ablation tests}

\paragraph{R1: \texttt{const\_0p5} (Sub~9).}
The information-theoretic floor: $-\ln 2 \approx -0.6931$. Anything
that does not beat this is worse than no model. We include it
because the gap between $-0.6931$ and our best (Sub~5 at
\lbNLLsubFive) quantifies how much of the leaderboard score comes
from prior calibration alone.

\paragraph{R2: \texttt{const\_glob\_mean} (Sub~10).}
The Bayes-optimal constant under squared error: $\bar y =
\sum_i y_i / N = 0.645$ from the training response matrix. If the
leaderboard's marginal pass rate matches training, this submission
recovers the entire R1\,$\to$\,R2 step from slide~29 ``for free.''

\paragraph{R3: \texttt{subj\_mean\_lookup} (Sub~15).}
The smoothed-subject-mean rung. We apply a 10\,\% Bayesian shrinkage
toward the global mean, $\hat p = 0.9\,\bar y_s + 0.1\,\bar
y_{\text{train}}$, because two of 745 training subjects
(\textit{DeepSeek-V3.2-Exp-Thinking} and
\textit{gemini-2.5-flash}) have $\bar y_s = 0$ exactly and a raw
lookup would risk an unbounded NLL contribution if those subjects
appear with $y{=}1$. Predictions are clipped to $[10^{-3}, 1{-}10^{-3}]$
per slide~28.

\paragraph{R4: \texttt{bm\_cond\_lookup} (Sub~16).}
Per-benchmark mean accuracy (we collapse the condition axis because
$>$80\,\% of items use \texttt{none}). This is expected to be the
weakest of the four prior submissions: leaderboard rounds may sample
held-out benchmarks not seen in training, in which case the
predictor degenerates to the global mean. The result therefore
isolates the contribution of \emph{generalising} metadata signals
versus memorising them.

\paragraph{R6 winner: \texttt{tscale\_mpnet} (Sub~5).}
Our best-performing submission. The all-mpnet-base-v2 encoder
produces 768-d embeddings for subject, item, benchmark, and
condition text; an MLP maps the concatenated representation to
per-item IRT parameters $(a, b)$; the full PGE
$\sigma(a\,(\theta - b))$ is computed in inference. Post-hoc
temperature scaling at $T{=}4.07$ (fit on a cold-start val split via
\texttt{dump\_cs\_logits.py}) flattens the over-confident logits
produced by the cold-start mismatch.

\paragraph{Sub~5 metadata blend: \texttt{sub5\_blend\_subj} (Sub~13).}
Slide~31 lists ``F4 metadata stabilisers'' as a hedge against the
content head being uncalibrated. We blend
$\hat p = 0.5\,p_{\text{sub5}} + 0.5\,\bar y_s$. The blend is
robust: T-scaled mpnet outputs are concentrated near $0.5$, so even
when $\bar y_s = 0$ the blend lands near $0.25$ rather than
saturating. This is the variant most likely to beat the Sub~5 NLL.

\paragraph{Sub~5 with K{=}5 adaptive labels: \texttt{sub5\_labeled\_shift} (Sub~14).}
Sub~8 attempted a per-subject mean-residual offset from the $K{=}5$
labels per category and regressed by $-0.10$ NLL—five labels split
across $\approx$5 subjects yields one label per subject, which is
much too noisy. We retreat to the most conservative possible
adaptive-label lever: a \emph{single} global offset
$\Delta = \frac{1}{N+\lambda}\sum_i (y_i - \hat p_i)$ with shrinkage
$\lambda{=}25$ (so 25 labels see a $25/(25{+}25) = 50\,\%$ pull) and a
magnitude cap $|\Delta| \le 0.05$. We intentionally omit
\texttt{labeling.py} so that the platform supplies uniformly random
labels, which give an unbiased mean-residual estimator.

\subsection{Strategies we did not pursue}
\label{sec:ablations-skipped}

\paragraph{F2 IRT alone (\ding{55}).}
Slide~31 warns this ``fails cold-start'': an IRT model without text
features cannot place a previously-unseen subject in the $\theta$
space. Our winning model already uses IRT \emph{with} text features
(rung~6); a clean ``IRT-only'' ablation would require retraining
from scratch on the train matrix without text, which we judged
unlikely to be informative.

\paragraph{F5 local LLM judge (\ding{55}).}
A small generative LLM scoring items on the sandbox GPU would
require a new HuggingFace repo declared in \texttt{models.txt}, a
calibration head trained on its log-probabilities, and a re-run of
the cold-start T-scaling sweep. The sandbox imposes a 30--60~min
per-tier timeout (Table~\ref{tab:gputiers}) and the implementation
budget did not permit a credible evaluation in the time remaining.

\subsection{Headline takeaway}

The pure-prior submissions (R1--R4) bound the contribution of
\emph{calibration alone}: the R1 floor is $-0.6931$ and Sub~5 sits at
\lbNLLsubFive. The remaining gap to the leaderboard ceiling
must therefore come from the content head and/or label-conditional
shifts. \texttt{sub5\_blend\_subj} (Sub~13) and
\texttt{sub5\_labeled\_shift} (Sub~14) target exactly those two
levers in the most conservative form we could justify against Sub~8's
$-0.10$ NLL regression.
` from
%  the main .tex, OR copy the body wholesale into the relevant section.
%
%  Required preamble packages (add to main.tex if not already there):
%      \usepackage{booktabs}
%      \usepackage{makecell}
%      \usepackage{array}
%      \usepackage{multirow}
%      \usepackage{xcolor}
%      \usepackage{amsmath}
%
%  Score macros — once the leaderboard returns each round, replace the
%  TBD with the number. Negative log-loss is signed so a sign is
%  required (e.g. -0.65, NOT 0.65).
% ====================================================================
\newcommand{\TBD}{\textit{pending}}
% -- known scores from prior rounds (sub 1..8) -----------------------
\newcommand{\lbNLLsubOne}{-1.01}\newcommand{\lbAUCsubOne}{0.62}
\newcommand{\lbNLLsubTwo}{-0.93}\newcommand{\lbAUCsubTwo}{0.63}
\newcommand{\lbNLLsubThree}{-0.85}\newcommand{\lbAUCsubThree}{0.63}
\newcommand{\lbNLLsubFive}{-0.65}\newcommand{\lbAUCsubFive}{0.63}
\newcommand{\lbNLLsubSix}{-0.65}\newcommand{\lbAUCsubSix}{\TBD}
\newcommand{\lbNLLsubSeven}{-0.65}\newcommand{\lbAUCsubSeven}{\TBD}
\newcommand{\lbNLLsubEight}{-0.75}\newcommand{\lbAUCsubEight}{\TBD}
% -- Part-V ablations: fill in as rounds resolve ---------------------
\newcommand{\lbNLLconstHalf}{\TBD}\newcommand{\lbAUCconstHalf}{\TBD}
\newcommand{\lbNLLconstGM}{\TBD}\newcommand{\lbAUCconstGM}{\TBD}
\newcommand{\lbNLLsubjMean}{\TBD}\newcommand{\lbAUCsubjMean}{\TBD}
\newcommand{\lbNLLbmCond}{\TBD}\newcommand{\lbAUCbmCond}{\TBD}
\newcommand{\lbNLLblendSubj}{\TBD}\newcommand{\lbAUCblendSubj}{\TBD}
\newcommand{\lbNLLlabeledShift}{\TBD}\newcommand{\lbAUClabeledShift}{\TBD}
% ====================================================================

\section{Ablations and Strategy Coverage}
\label{sec:ablations}

We organise our exploration around Discussion~6's two scaffolds:
(i)~the seven-rung calibration climb (from a constant 0.5 up to a
content-aware PGE with adaptive labels), and
(ii)~the six-family option space
(F1 calibrated priors, F2 IRT alone, F3 content-based, F4 metadata
stabilisers, F5 local LLM judges, F6 adaptive labels).
Table~\ref{tab:coverage} maps each submission we ran onto these axes,
and Table~\ref{tab:ablations} reports the leaderboard outcomes.

\subsection{Coverage map}

\begin{table}[h]
\centering
\small
\caption{Coverage of Discussion 6 strategy axes. Rows are the
seven-rung calibration climb; columns are the six families.
\ding{51} marks an ablation we shipped to the leaderboard; \ding{55} marks
families we deliberately did not pursue (with justification in
\S\ref{sec:ablations-skipped}).}
\label{tab:coverage}
\begin{tabular}{lcccccc}
\toprule
                                  & \textbf{F1} & \textbf{F2} & \textbf{F3} & \textbf{F4} & \textbf{F5} & \textbf{F6} \\
                                  & priors      & IRT         & content     & metadata    & local LLM   & adaptive \\
\midrule
R1: $\hat p = 0.5$                & \checkmark  &             &             &             &             &             \\
R2: global mean                   & \checkmark  &             &             &             &             &             \\
R3: smoothed subject mean         & \checkmark  &             &             & \checkmark  &             &             \\
R4: per-(benchmark, condition) mean & \checkmark &             &             & \checkmark  &             &             \\
R5: + content (text $\to$ params) &             &             & \checkmark  &             & \ding{55}   &             \\
R6: + IRT $\theta$ (full PGE)     &             & \ding{55}   & \checkmark  &             &             &             \\
R7: + adaptive K{=}5 labels       &             &             & \checkmark  &             &             & \checkmark  \\
\bottomrule
\end{tabular}
\end{table}

The lower four rungs are pure ablations of slide~30's recipe; the
upper three are variants of our winning submission (Sub~5, T-scaled
mpnet) along the metadata and adaptive-label axes.

\subsection{Results}

\begin{table*}[ht]
\centering
\small
\caption{All submissions, ordered by leaderboard NLL (negative
log-loss; higher is better). LB columns are filled in as each round
resolves. Rungs and families reference Discussion~6 Part~V
slides~29 and~31.}
\label{tab:ablations}
\renewcommand{\arraystretch}{1.15}
\begin{tabular}{rl c c l rr}
\toprule
\textbf{ID} & \textbf{Submission} & \textbf{Rung} & \textbf{Family} & \textbf{Mechanism} & \textbf{LB NLL} & \textbf{LB AUC} \\
\midrule
\multicolumn{7}{l}{\textit{Climb baselines (priors / metadata)}}\\
9  & \texttt{const\_0p5}       & R1 & F1     & constant $0.5$                                       & \lbNLLconstHalf      & \lbAUCconstHalf \\
10 & \texttt{const\_glob\_mean} & R2 & F1     & constant $\bar y_{\text{train}}{=}0.645$             & \lbNLLconstGM        & \lbAUCconstGM   \\
15 & \texttt{subj\_mean\_lookup}& R3 & F1+F4  & $0.9\,\bar y_s + 0.1\,\bar y_{\text{train}}$         & \lbNLLsubjMean       & \lbAUCsubjMean  \\
16 & \texttt{bm\_cond\_lookup} & R4 & F1+F4  & $\bar y_{(b,c)}$, fallback $\bar y_{\text{train}}$   & \lbNLLbmCond         & \lbAUCbmCond    \\
\midrule
\multicolumn{7}{l}{\textit{Prior context: encoder upgrade alone vs. with T-scaling}}\\
1 & \texttt{mpnet\_platt\_meanresid}     & R6 & F3 & mpnet+Platt+per-subject offset  & \lbNLLsubOne   & \lbAUCsubOne   \\
2 & \texttt{mpnet\_noplatt\_meanresid}   & R6 & F3 & sub 1 -- Platt                  & \lbNLLsubTwo   & \lbAUCsubTwo   \\
3 & \texttt{mpnet\_no\_offset}           & R6 & F3 & sub 2 -- offset                 & \lbNLLsubThree & \lbAUCsubThree \\
\textbf{5} & \textbf{\texttt{tscale\_mpnet} (winner)} & R6 & F3 & \textbf{sub 3 + post-hoc T-scaling ($T{=}4.07$)} & \textbf{\lbNLLsubFive} & \textbf{\lbAUCsubFive} \\
6  & \texttt{tscale\_T3p0}              & R6 & F3 & sub 5 with $T{=}3.0$              & \lbNLLsubSix    & \lbAUCsubSix     \\
7  & \texttt{tscale\_T5p5}              & R6 & F3 & sub 5 with $T{=}5.5$              & \lbNLLsubSeven  & \lbAUCsubSeven   \\
8  & \texttt{tscale\_plus\_offset}      & R6 & F3+F4 & sub 5 + per-subject mean-resid    & \lbNLLsubEight  & \lbAUCsubEight   \\
\midrule
\multicolumn{7}{l}{\textit{Part-V ablations on top of Sub 5}}\\
13 & \texttt{sub5\_blend\_subj}     & R6+R3 & F3+F4 & $0.5\,p_{\text{sub5}} + 0.5\,\bar y_s$       & \lbNLLblendSubj    & \lbAUCblendSubj    \\
14 & \texttt{sub5\_labeled\_shift}  & R7    & F3+F6 & sub 5 + shrunk global offset from $K{=}5$ labels, $\lambda{=}25$, $|\Delta|{\le}0.05$ & \lbNLLlabeledShift & \lbAUClabeledShift \\
\bottomrule
\end{tabular}
\end{table*}

\subsection{What each ablation tests}

\paragraph{R1: \texttt{const\_0p5} (Sub~9).}
The information-theoretic floor: $-\ln 2 \approx -0.6931$. Anything
that does not beat this is worse than no model. We include it
because the gap between $-0.6931$ and our best (Sub~5 at
\lbNLLsubFive) quantifies how much of the leaderboard score comes
from prior calibration alone.

\paragraph{R2: \texttt{const\_glob\_mean} (Sub~10).}
The Bayes-optimal constant under squared error: $\bar y =
\sum_i y_i / N = 0.645$ from the training response matrix. If the
leaderboard's marginal pass rate matches training, this submission
recovers the entire R1\,$\to$\,R2 step from slide~29 ``for free.''

\paragraph{R3: \texttt{subj\_mean\_lookup} (Sub~15).}
The smoothed-subject-mean rung. We apply a 10\,\% Bayesian shrinkage
toward the global mean, $\hat p = 0.9\,\bar y_s + 0.1\,\bar
y_{\text{train}}$, because two of 745 training subjects
(\textit{DeepSeek-V3.2-Exp-Thinking} and
\textit{gemini-2.5-flash}) have $\bar y_s = 0$ exactly and a raw
lookup would risk an unbounded NLL contribution if those subjects
appear with $y{=}1$. Predictions are clipped to $[10^{-3}, 1{-}10^{-3}]$
per slide~28.

\paragraph{R4: \texttt{bm\_cond\_lookup} (Sub~16).}
Per-benchmark mean accuracy (we collapse the condition axis because
$>$80\,\% of items use \texttt{none}). This is expected to be the
weakest of the four prior submissions: leaderboard rounds may sample
held-out benchmarks not seen in training, in which case the
predictor degenerates to the global mean. The result therefore
isolates the contribution of \emph{generalising} metadata signals
versus memorising them.

\paragraph{R6 winner: \texttt{tscale\_mpnet} (Sub~5).}
Our best-performing submission. The all-mpnet-base-v2 encoder
produces 768-d embeddings for subject, item, benchmark, and
condition text; an MLP maps the concatenated representation to
per-item IRT parameters $(a, b)$; the full PGE
$\sigma(a\,(\theta - b))$ is computed in inference. Post-hoc
temperature scaling at $T{=}4.07$ (fit on a cold-start val split via
\texttt{dump\_cs\_logits.py}) flattens the over-confident logits
produced by the cold-start mismatch.

\paragraph{Sub~5 metadata blend: \texttt{sub5\_blend\_subj} (Sub~13).}
Slide~31 lists ``F4 metadata stabilisers'' as a hedge against the
content head being uncalibrated. We blend
$\hat p = 0.5\,p_{\text{sub5}} + 0.5\,\bar y_s$. The blend is
robust: T-scaled mpnet outputs are concentrated near $0.5$, so even
when $\bar y_s = 0$ the blend lands near $0.25$ rather than
saturating. This is the variant most likely to beat the Sub~5 NLL.

\paragraph{Sub~5 with K{=}5 adaptive labels: \texttt{sub5\_labeled\_shift} (Sub~14).}
Sub~8 attempted a per-subject mean-residual offset from the $K{=}5$
labels per category and regressed by $-0.10$ NLL—five labels split
across $\approx$5 subjects yields one label per subject, which is
much too noisy. We retreat to the most conservative possible
adaptive-label lever: a \emph{single} global offset
$\Delta = \frac{1}{N+\lambda}\sum_i (y_i - \hat p_i)$ with shrinkage
$\lambda{=}25$ (so 25 labels see a $25/(25{+}25) = 50\,\%$ pull) and a
magnitude cap $|\Delta| \le 0.05$. We intentionally omit
\texttt{labeling.py} so that the platform supplies uniformly random
labels, which give an unbiased mean-residual estimator.

\subsection{Strategies we did not pursue}
\label{sec:ablations-skipped}

\paragraph{F2 IRT alone (\ding{55}).}
Slide~31 warns this ``fails cold-start'': an IRT model without text
features cannot place a previously-unseen subject in the $\theta$
space. Our winning model already uses IRT \emph{with} text features
(rung~6); a clean ``IRT-only'' ablation would require retraining
from scratch on the train matrix without text, which we judged
unlikely to be informative.

\paragraph{F5 local LLM judge (\ding{55}).}
A small generative LLM scoring items on the sandbox GPU would
require a new HuggingFace repo declared in \texttt{models.txt}, a
calibration head trained on its log-probabilities, and a re-run of
the cold-start T-scaling sweep. The sandbox imposes a 30--60~min
per-tier timeout (Table~\ref{tab:gputiers}) and the implementation
budget did not permit a credible evaluation in the time remaining.

\subsection{Headline takeaway}

The pure-prior submissions (R1--R4) bound the contribution of
\emph{calibration alone}: the R1 floor is $-0.6931$ and Sub~5 sits at
\lbNLLsubFive. The remaining gap to the leaderboard ceiling
must therefore come from the content head and/or label-conditional
shifts. \texttt{sub5\_blend\_subj} (Sub~13) and
\texttt{sub5\_labeled\_shift} (Sub~14) target exactly those two
levers in the most conservative form we could justify against Sub~8's
$-0.10$ NLL regression.
` from
%  the main .tex, OR copy the body wholesale into the relevant section.
%
%  Required preamble packages (add to main.tex if not already there):
%      \usepackage{booktabs}
%      \usepackage{makecell}
%      \usepackage{array}
%      \usepackage{multirow}
%      \usepackage{xcolor}
%      \usepackage{amsmath}
%
%  Score macros — once the leaderboard returns each round, replace the
%  TBD with the number. Negative log-loss is signed so a sign is
%  required (e.g. -0.65, NOT 0.65).
% ====================================================================
\newcommand{\TBD}{\textit{pending}}
% -- known scores from prior rounds (sub 1..8) -----------------------
\newcommand{\lbNLLsubOne}{-1.01}\newcommand{\lbAUCsubOne}{0.62}
\newcommand{\lbNLLsubTwo}{-0.93}\newcommand{\lbAUCsubTwo}{0.63}
\newcommand{\lbNLLsubThree}{-0.85}\newcommand{\lbAUCsubThree}{0.63}
\newcommand{\lbNLLsubFive}{-0.65}\newcommand{\lbAUCsubFive}{0.63}
\newcommand{\lbNLLsubSix}{-0.65}\newcommand{\lbAUCsubSix}{\TBD}
\newcommand{\lbNLLsubSeven}{-0.65}\newcommand{\lbAUCsubSeven}{\TBD}
\newcommand{\lbNLLsubEight}{-0.75}\newcommand{\lbAUCsubEight}{\TBD}
% -- Part-V ablations: fill in as rounds resolve ---------------------
\newcommand{\lbNLLconstHalf}{\TBD}\newcommand{\lbAUCconstHalf}{\TBD}
\newcommand{\lbNLLconstGM}{\TBD}\newcommand{\lbAUCconstGM}{\TBD}
\newcommand{\lbNLLsubjMean}{\TBD}\newcommand{\lbAUCsubjMean}{\TBD}
\newcommand{\lbNLLbmCond}{\TBD}\newcommand{\lbAUCbmCond}{\TBD}
\newcommand{\lbNLLblendSubj}{\TBD}\newcommand{\lbAUCblendSubj}{\TBD}
\newcommand{\lbNLLlabeledShift}{\TBD}\newcommand{\lbAUClabeledShift}{\TBD}
% ====================================================================

\section{Ablations and Strategy Coverage}
\label{sec:ablations}

We organise our exploration around Discussion~6's two scaffolds:
(i)~the seven-rung calibration climb (from a constant 0.5 up to a
content-aware PGE with adaptive labels), and
(ii)~the six-family option space
(F1 calibrated priors, F2 IRT alone, F3 content-based, F4 metadata
stabilisers, F5 local LLM judges, F6 adaptive labels).
Table~\ref{tab:coverage} maps each submission we ran onto these axes,
and Table~\ref{tab:ablations} reports the leaderboard outcomes.

\subsection{Coverage map}

\begin{table}[h]
\centering
\small
\caption{Coverage of Discussion 6 strategy axes. Rows are the
seven-rung calibration climb; columns are the six families.
\ding{51} marks an ablation we shipped to the leaderboard; \ding{55} marks
families we deliberately did not pursue (with justification in
\S\ref{sec:ablations-skipped}).}
\label{tab:coverage}
\begin{tabular}{lcccccc}
\toprule
                                  & \textbf{F1} & \textbf{F2} & \textbf{F3} & \textbf{F4} & \textbf{F5} & \textbf{F6} \\
                                  & priors      & IRT         & content     & metadata    & local LLM   & adaptive \\
\midrule
R1: $\hat p = 0.5$                & \checkmark  &             &             &             &             &             \\
R2: global mean                   & \checkmark  &             &             &             &             &             \\
R3: smoothed subject mean         & \checkmark  &             &             & \checkmark  &             &             \\
R4: per-(benchmark, condition) mean & \checkmark &             &             & \checkmark  &             &             \\
R5: + content (text $\to$ params) &             &             & \checkmark  &             & \ding{55}   &             \\
R6: + IRT $\theta$ (full PGE)     &             & \ding{55}   & \checkmark  &             &             &             \\
R7: + adaptive K{=}5 labels       &             &             & \checkmark  &             &             & \checkmark  \\
\bottomrule
\end{tabular}
\end{table}

The lower four rungs are pure ablations of slide~30's recipe; the
upper three are variants of our winning submission (Sub~5, T-scaled
mpnet) along the metadata and adaptive-label axes.

\subsection{Results}

\begin{table*}[ht]
\centering
\small
\caption{All submissions, ordered by leaderboard NLL (negative
log-loss; higher is better). LB columns are filled in as each round
resolves. Rungs and families reference Discussion~6 Part~V
slides~29 and~31.}
\label{tab:ablations}
\renewcommand{\arraystretch}{1.15}
\begin{tabular}{rl c c l rr}
\toprule
\textbf{ID} & \textbf{Submission} & \textbf{Rung} & \textbf{Family} & \textbf{Mechanism} & \textbf{LB NLL} & \textbf{LB AUC} \\
\midrule
\multicolumn{7}{l}{\textit{Climb baselines (priors / metadata)}}\\
9  & \texttt{const\_0p5}       & R1 & F1     & constant $0.5$                                       & \lbNLLconstHalf      & \lbAUCconstHalf \\
10 & \texttt{const\_glob\_mean} & R2 & F1     & constant $\bar y_{\text{train}}{=}0.645$             & \lbNLLconstGM        & \lbAUCconstGM   \\
15 & \texttt{subj\_mean\_lookup}& R3 & F1+F4  & $0.9\,\bar y_s + 0.1\,\bar y_{\text{train}}$         & \lbNLLsubjMean       & \lbAUCsubjMean  \\
16 & \texttt{bm\_cond\_lookup} & R4 & F1+F4  & $\bar y_{(b,c)}$, fallback $\bar y_{\text{train}}$   & \lbNLLbmCond         & \lbAUCbmCond    \\
\midrule
\multicolumn{7}{l}{\textit{Prior context: encoder upgrade alone vs. with T-scaling}}\\
1 & \texttt{mpnet\_platt\_meanresid}     & R6 & F3 & mpnet+Platt+per-subject offset  & \lbNLLsubOne   & \lbAUCsubOne   \\
2 & \texttt{mpnet\_noplatt\_meanresid}   & R6 & F3 & sub 1 -- Platt                  & \lbNLLsubTwo   & \lbAUCsubTwo   \\
3 & \texttt{mpnet\_no\_offset}           & R6 & F3 & sub 2 -- offset                 & \lbNLLsubThree & \lbAUCsubThree \\
\textbf{5} & \textbf{\texttt{tscale\_mpnet} (winner)} & R6 & F3 & \textbf{sub 3 + post-hoc T-scaling ($T{=}4.07$)} & \textbf{\lbNLLsubFive} & \textbf{\lbAUCsubFive} \\
6  & \texttt{tscale\_T3p0}              & R6 & F3 & sub 5 with $T{=}3.0$              & \lbNLLsubSix    & \lbAUCsubSix     \\
7  & \texttt{tscale\_T5p5}              & R6 & F3 & sub 5 with $T{=}5.5$              & \lbNLLsubSeven  & \lbAUCsubSeven   \\
8  & \texttt{tscale\_plus\_offset}      & R6 & F3+F4 & sub 5 + per-subject mean-resid    & \lbNLLsubEight  & \lbAUCsubEight   \\
\midrule
\multicolumn{7}{l}{\textit{Part-V ablations on top of Sub 5}}\\
13 & \texttt{sub5\_blend\_subj}     & R6+R3 & F3+F4 & $0.5\,p_{\text{sub5}} + 0.5\,\bar y_s$       & \lbNLLblendSubj    & \lbAUCblendSubj    \\
14 & \texttt{sub5\_labeled\_shift}  & R7    & F3+F6 & sub 5 + shrunk global offset from $K{=}5$ labels, $\lambda{=}25$, $|\Delta|{\le}0.05$ & \lbNLLlabeledShift & \lbAUClabeledShift \\
\bottomrule
\end{tabular}
\end{table*}

\subsection{What each ablation tests}

\paragraph{R1: \texttt{const\_0p5} (Sub~9).}
The information-theoretic floor: $-\ln 2 \approx -0.6931$. Anything
that does not beat this is worse than no model. We include it
because the gap between $-0.6931$ and our best (Sub~5 at
\lbNLLsubFive) quantifies how much of the leaderboard score comes
from prior calibration alone.

\paragraph{R2: \texttt{const\_glob\_mean} (Sub~10).}
The Bayes-optimal constant under squared error: $\bar y =
\sum_i y_i / N = 0.645$ from the training response matrix. If the
leaderboard's marginal pass rate matches training, this submission
recovers the entire R1\,$\to$\,R2 step from slide~29 ``for free.''

\paragraph{R3: \texttt{subj\_mean\_lookup} (Sub~15).}
The smoothed-subject-mean rung. We apply a 10\,\% Bayesian shrinkage
toward the global mean, $\hat p = 0.9\,\bar y_s + 0.1\,\bar
y_{\text{train}}$, because two of 745 training subjects
(\textit{DeepSeek-V3.2-Exp-Thinking} and
\textit{gemini-2.5-flash}) have $\bar y_s = 0$ exactly and a raw
lookup would risk an unbounded NLL contribution if those subjects
appear with $y{=}1$. Predictions are clipped to $[10^{-3}, 1{-}10^{-3}]$
per slide~28.

\paragraph{R4: \texttt{bm\_cond\_lookup} (Sub~16).}
Per-benchmark mean accuracy (we collapse the condition axis because
$>$80\,\% of items use \texttt{none}). This is expected to be the
weakest of the four prior submissions: leaderboard rounds may sample
held-out benchmarks not seen in training, in which case the
predictor degenerates to the global mean. The result therefore
isolates the contribution of \emph{generalising} metadata signals
versus memorising them.

\paragraph{R6 winner: \texttt{tscale\_mpnet} (Sub~5).}
Our best-performing submission. The all-mpnet-base-v2 encoder
produces 768-d embeddings for subject, item, benchmark, and
condition text; an MLP maps the concatenated representation to
per-item IRT parameters $(a, b)$; the full PGE
$\sigma(a\,(\theta - b))$ is computed in inference. Post-hoc
temperature scaling at $T{=}4.07$ (fit on a cold-start val split via
\texttt{dump\_cs\_logits.py}) flattens the over-confident logits
produced by the cold-start mismatch.

\paragraph{Sub~5 metadata blend: \texttt{sub5\_blend\_subj} (Sub~13).}
Slide~31 lists ``F4 metadata stabilisers'' as a hedge against the
content head being uncalibrated. We blend
$\hat p = 0.5\,p_{\text{sub5}} + 0.5\,\bar y_s$. The blend is
robust: T-scaled mpnet outputs are concentrated near $0.5$, so even
when $\bar y_s = 0$ the blend lands near $0.25$ rather than
saturating. This is the variant most likely to beat the Sub~5 NLL.

\paragraph{Sub~5 with K{=}5 adaptive labels: \texttt{sub5\_labeled\_shift} (Sub~14).}
Sub~8 attempted a per-subject mean-residual offset from the $K{=}5$
labels per category and regressed by $-0.10$ NLL—five labels split
across $\approx$5 subjects yields one label per subject, which is
much too noisy. We retreat to the most conservative possible
adaptive-label lever: a \emph{single} global offset
$\Delta = \frac{1}{N+\lambda}\sum_i (y_i - \hat p_i)$ with shrinkage
$\lambda{=}25$ (so 25 labels see a $25/(25{+}25) = 50\,\%$ pull) and a
magnitude cap $|\Delta| \le 0.05$. We intentionally omit
\texttt{labeling.py} so that the platform supplies uniformly random
labels, which give an unbiased mean-residual estimator.

\subsection{Strategies we did not pursue}
\label{sec:ablations-skipped}

\paragraph{F2 IRT alone (\ding{55}).}
Slide~31 warns this ``fails cold-start'': an IRT model without text
features cannot place a previously-unseen subject in the $\theta$
space. Our winning model already uses IRT \emph{with} text features
(rung~6); a clean ``IRT-only'' ablation would require retraining
from scratch on the train matrix without text, which we judged
unlikely to be informative.

\paragraph{F5 local LLM judge (\ding{55}).}
A small generative LLM scoring items on the sandbox GPU would
require a new HuggingFace repo declared in \texttt{models.txt}, a
calibration head trained on its log-probabilities, and a re-run of
the cold-start T-scaling sweep. The sandbox imposes a 30--60~min
per-tier timeout (Table~\ref{tab:gputiers}) and the implementation
budget did not permit a credible evaluation in the time remaining.

\subsection{Headline takeaway}

The pure-prior submissions (R1--R4) bound the contribution of
\emph{calibration alone}: the R1 floor is $-0.6931$ and Sub~5 sits at
\lbNLLsubFive. The remaining gap to the leaderboard ceiling
must therefore come from the content head and/or label-conditional
shifts. \texttt{sub5\_blend\_subj} (Sub~13) and
\texttt{sub5\_labeled\_shift} (Sub~14) target exactly those two
levers in the most conservative form we could justify against Sub~8's
$-0.10$ NLL regression.
